% Part: first-order-logic
% Chapter: syntax-and-semantics
% Section: subformulas

\documentclass[../../../include/open-logic-section]{subfiles}

\begin{document}

\olfileid{fol}{syn}{sbf}

\olsection{\printtoken{P}{subformula}}

\begin{explain}
讨论“组成”给定公式的那些公式往往很有用。
我们称它们为该公式的\emph{子公式}。
任何公式都算作它自身的子公式；$!A$ 的子公式中除 $!A$ 自身以外的称为\emph{真子公式}。
\end{explain}

\begin{defn}[直接子公式]
如果 $!A$ 是一个公式，则 $!A$ 的\emph{直接子公式}归纳定义如下：
\begin{enumerate}
\item 原子公式没有直接子公式。

\tagitem{prvNot}{\indcase{!A}{\lnot !B}{$\indfrm$ 的唯一直接子公式是~$!B$。}}{}

\item \indcase{!A}{(!B \ast !C)}{$\indfrm$ 的直接子公式是 $!B$ 和 $!C$（$\ast$ 是任一二元联结词）。}

\tagitem{prvAll}{\indcase{!A}{\lforall[x][!B]}{$\indfrm$ 的唯一直接子公式是~$!B$。}}{}

\tagitem{prvEx}{\indcase{!A}{\lexists[x][!B]}{$\indfrm$ 的唯一直接子公式是~$!B$。}}{}
\end{enumerate}
\end{defn}

\begin{defn}[真子公式]
如果 $!A$ 是一个公式，则 $!A$ 的\emph{真子公式}递归定义如下：
\begin{enumerate}
\item 原子公式没有真子公式。

\tagitem{prvNot}{\indcase{!A}{\lnot !B}{$\indfrm$ 的真子公式是 $!B$ 以及 $!B$ 的所有真子公式。}}{}

\item \indcase{!A}{(!B \ast !C)}{$\indfrm$ 的真子公式是 $!B$、$!C$，以及 $!B$ 的所有真子公式和 $!C$ 的所有真子公式。}

\tagitem{prvAll}{\indcase{!A}{\lforall[x][!B]}{$\indfrm$ 的真子公式是 $!B$ 以及 $!B$ 的所有真子公式。}}{}

\tagitem{prvEx}{\indcase{!A}{\lexists[x][!B]}{$\indfrm$ 的真子公式是 $!B$ 以及 $!B$ 的所有真子公式。}}{}
\end{enumerate}
\end{defn}

\begin{defn}[子公式]
$!A$ 的\emph{子公式}是 $!A$ 自身及其所有真子公式。
\end{defn}

\begin{explain}
请注意我们在定义直接子公式与真子公式时的微妙差别。
对于第一种情况，我们直接针对公式~$!A$ 的每种可能形式，定义了它的直接子公式。
这是一种分情况显式定义，其分情况的方式与公式集合的归纳定义相照应。
对于第二种情况，我们也照应了所有公式集合的定义方式，
但在每种情况中，除了较小的公式 $!B$、$!C$ 自身之外，我们还包含了它们的真子公式。
这使得该定义成为\emph{递归的}。
一般而言，在归纳定义的集合（此处为公式）上定义函数时，如果函数定义的各分支用到了函数本身，那么这个定义就是递归的。
然而，为确保良定义性，我们必须保证，只使用函数在“先于”当前所定义的自变元上的值——在我们的例子中，当定义 $(!B \ast !C)$ 的“真子公式”时，我们只使用了“更早”公式 $!B$ 和 $!C$ 的真子公式。
\end{explain}

\begin{prop}
\ollabel{prop:subfrm-trans}
假设 $!B$ 是 $!A$ 的子公式且 $!C$ 是 $!B$ 的子公式，
那么 $!C$ 也是 $!A$ 的子公式。换言之，子公式关系是传递的。
\end{prop}

\begin{prob}
证明 \olref[fol][syn][sbf]{prop:subfrm-trans}。
\end{prob}

\begin{prop}
\ollabel{prop:count-subfrms}
假设公式 $!A$ 含有 $n$ 个联结词和量词，
那么 $!A$ 至多有 $2n+1$ 个子公式。
\end{prop}

\begin{prob}
证明 \olref[fol][syn][sbf]{prop:count-subfrms}。
\end{prob}

\end{document}